% Source: https://artofproblemsolving.com/community/c3h1673126_circles_xd
% Stable identity: aops-topic-1673126-q01
Let the centres be $O_1$, $O_2$ and the intersection points $P$, $Q$. Then $O_1O_2=O_1P=O_2P=r$, so triangles $O_1O_2P$ and $O_1O_2Q$ are equilateral, and $\angle PO_1Q=\angle PO_2Q=120^\circ$.

\begin{asy}
usepackage("amsmath");
size(5.5cm);
defaultpen(fontsize(10pt));
real r=2;pair O1=(0,0),O2=(r,0),P=(r/2,r*sqrt(3)/2),Q=(r/2,-r*sqrt(3)/2);
fill(arc(O1,r,-60,60)--arc(O2,r,120,240)--cycle,paleblue);
draw(circle(O1,r),linewidth(0.9));draw(circle(O2,r),linewidth(0.9));
draw(O1--P--O2--Q--cycle,blue+dashed);draw(O1--O2,blue+dashed);draw(P--Q,dashed);
dot(O1);dot(O2);dot(P);dot(Q);
label("$O_1$",O1,W);label("$O_2$",O2,E);label("$P$",P,N);label("$Q$",Q,S);
label("$r$",(O1+O2)/4,S);
label("\mbox{}",(min(currentpicture,user=true)+max(currentpicture,user=true))/2,black); // keeps the text after the figure black when LaTeX typesets the labels inline
\end{asy}

\textbf{(a)} The common region (shaded) is cut by the chord $PQ$ into two congruent circular segments, each equal to the $120^\circ$ sector $O_1PQ$ minus triangle $O_1PQ$:
\[
\dfrac{120}{360}\pi r^2-\dfrac{1}{2}r^2\sin120^\circ=\dfrac{\pi r^2}{3}-\dfrac{\sqrt3}{4}r^2 .
\]
So the area is $\boxed{\left(\dfrac{2\pi}{3}-\dfrac{\sqrt3}{2}\right)r^2}\approx1.228r^2$. Its boundary consists of two $120^\circ$ arcs, so its perimeter is $2\cdot\dfrac{1}{3}\cdot2\pi r=\boxed{\dfrac{4\pi r}{3}}$.

\textbf{(b)} By inclusion-exclusion the area of the union is
\[
2\pi r^2-\left(\dfrac{2\pi}{3}-\dfrac{\sqrt3}{2}\right)r^2=\boxed{\left(\dfrac{4\pi}{3}+\dfrac{\sqrt3}{2}\right)r^2},
\]
and its boundary consists of the two $240^\circ$ arcs outside the other circle, of total length $2\cdot\dfrac{2}{3}\cdot2\pi r=\boxed{\dfrac{8\pi r}{3}}$.
